\documentclass[a4paper]{article}

\usepackage{graphicx}
\usepackage{amsmath,amssymb}
\usepackage{minted}
\usepackage{multirow}
\usepackage{float}
\usepackage{todonotes}
\usepackage{forest}
\usepackage{xstring}
\usepackage{xcolor}
\usepackage[most]{tcolorbox}
\usepackage{xparse}
\usepackage{natbib}

\usetikzlibrary{graphs,graphdrawing}
\usegdlibrary{layered}

% for external graphviz
\input{/home/donkarlo/Dropbox/repo/nd_language_ai_project/src/nd_language_ai/formal/computer/programming/tex/latex/code_clips/external_graphviz.tex}

% for tags
\input{/home/donkarlo/Dropbox/repo/nd_language_ai_project/src/nd_language_ai/formal/computer/programming/tex/latex/parts/control_sequence/kind/semtag/semtag.tex}

% for links
\input{/home/donkarlo/Dropbox/repo/nd_language_ai_project/src/nd_language_ai/formal/computer/programming/tex/latex/code_clips/seehere.tex}

\title{Kinematic teaching}
\date{}
\author{Mohammad Rahmani}

\begin{document}
    \maketitle
    The paper uses kinesthetic teaching to collect demonstrations for contact-rich robotic manipulation.
    Visual and tactile information are jointly used to learn manipulation policies from human demonstrations.
    Tactile force matching modifies demonstrated trajectories so that the robot reproduces the forces observed during teaching. Learned mode switching allows the policy to adapt between free-space motion and physical contact using visual and tactile feedback.Experiments on door-opening tasks show that force matching and visuotactile sensing substantially improve imitation-learning success rates \cite{ablett-2025-multimodal-and-force-matched-imitation-learning-with-a-see-through-visuotactile-sensor}.
    \bibliographystyle{plainnat}
    \bibliography{/home/donkarlo/Dropbox/repo/nd_shared_research/bibliography/bibliography.bib}
\end{document}